\documentclass[twocolumn]{aastex631}
\usepackage{amsmath}
\usepackage{booktabs}
\shorttitle{SDSS density proxy for environmental quenching}
\shortauthors{NebulaMind local integration}
\begin{document}

\title{SDSS density proxy for environmental quenching: selection-aware SDSS optical proxy integration}
\author{NebulaMind Research Autopilot}
\affiliation{Local reproducible integration run; public SDSS DR17 data only}

\begin{abstract}
We integrate the active proposal 'Separating internal and environmental quenching across stellar mass, halo mass, and redshift' as a guarded SDSS DR17 optical denominator/proxy draft rather than as a completed physical-feedback paper. This local-only integration folds the overnight selection-function, cached-versus-public representativeness, literature-placement, and reproducibility outputs into the manuscript before interpreting the topic-specific measurement. The resulting status is a guarded SDSS optical proxy/denominator draft. No public page, live root, database, deployment, git, or external submission action is part of this run.
\end{abstract}

\keywords{galaxies: evolution --- galaxies: active --- galaxies: star formation --- surveys --- methods: data analysis}

\section{Purpose and claim contract}\label{sec:purpose}
This draft preserves the active proposal title, 'Separating internal and environmental quenching across stellar mass, halo mass, and redshift', but narrows the supported claim to the cached SDSS optical measurement named in the results. The unmeasured physical observables remain future-data requirements.

The claim contract is intentionally conservative. Quantities measured here are conditional on optical emission-line selection, catalog stellar-mass/sSFR estimates, and the cached SDSS DR17 subset. Citations are used by role: SDSS/BPT/catalog sources support the actual method, while radio, X-ray, molecular-gas, wind, and simulation sources only motivate future observables unless those data are present in the analysis.


\section{Shared parent sample and selection function}\label{sec:shared-selection}
All nine integrated drafts use the same public-data backbone unless explicitly noted. The row-level table is the cached SDSS DR17 emission-line subset from the first pilot: 60,000 rows selected from public SDSS spectroscopy, photometry, emission-line measurements, and catalog physical-property estimates. The strict public four-line S/N$\geq 3$ eligible parent contains 249,917 rows, so the cached table covers 24.0\% of that strict parent. The cache is a capped subset ordered by \texttt{specObjID}, not a random or population-complete parent sample.

\begin{deluxetable*}{lrrr}
\tabletypesize{\scriptsize}
\tablecaption{Shared SDSS DR17 selection cascade used before paper-specific quantities.\label{tab:selection-cascade}}
\tablehead{\colhead{Selection stage} & \colhead{Public DR17 rows} & \colhead{Cached rows} & \colhead{Retention vs. spectro-z parent}}
\startdata
SpecObj GALAXY, 0.02<z<0.12 & 501,060 & -- & 1.000 \\
plus galSpecInfo/PhotoObj/galSpecExtra and mass/sSFR bounds & 416,554 & -- & 0.831 \\
plus galSpecLine join & 416,554 & -- & 0.831 \\
four BPT lines positive with positive errors & 373,445 & 60,000 & 0.745 \\
four BPT lines S/N>=3 & 249,917 & 60,000 & 0.499 \\
four BPT lines S/N>=5 & 176,523 & 42,446 & 0.352 \\
four BPT lines S/N>=10 & 91,768 & 22,311 & 0.183 \\
\enddata
\tablecomments{Counts are read-only public SDSS DR17 count queries plus the cached local CSV. Cached rows are shown only where the cache applies.}
\end{deluxetable*}

The four-line requirement is strongly selection dependent. In the public counts, S/N$\geq3$ keeps 33.6\% of the $-12<\log {\rm sSFR}<-11$ parent bin but 94.9\% of the $-10<\log {\rm sSFR}<-9.5$ bin. Therefore every incidence, quenching, density, gas-denominator, or target-vector statement in this integration is conditional on the four-line emission-line selection.

Cached-versus-public marginal checks found no redshift, stellar-mass, or sSFR bin with a cached-minus-public fraction difference above 5 percentage points. The largest absolute differences were 2.03 percentage points in redshift, -1.63 percentage points in stellar mass, and -0.58 percentage points in sSFR. This is a representativeness diagnostic only; it does not make the capped cache random or complete.


\section{Measurements}\label{sec:measurements}
The row-level measurements include redshift, stellar mass, catalog specific star-formation rate, model colors, and four optical line fluxes/errors. BPT labels and all derived denominators are recomputed locally from the cached table. Catalog SFR/sSFR values are treated as low-redshift SDSS physical-property estimates rather than direct resolved gas or feedback measurements \citep{sdssdr17,brinchmann2004,york2000}.


\section{Topic-specific optical denominator or proxy result}\label{sec:topic-result}
The consolidated proposal question is: Does a nearest-neighbour density proxy add quenched-fraction information beyond stellar mass in the SDSS emission-line sample? The integrated manuscript deliberately demotes the claim to the directly measured SDSS quantity. It is not a full physical-feedback test.

\begin{itemize}
\item The SDSS emission-line denominator contains 60,000 galaxies with an internally computed 10th-neighbour density proxy.
\item The high-density quartile has quenched fraction 0.230 (3,456/15,000); the low-density quartile has 0.181 (2,710/15,000).
\item The bootstrap high-minus-low quenched-fraction interval is [0.041, 0.059].
\item A linear probability model adjusted for log stellar mass and redshift gives a high-density coefficient of 0.032 +/- 0.004.
\end{itemize}


\begin{figure}
\centering
\includegraphics[width=\columnwidth]{../figures/fig-topic.pdf}
\caption{Topic-specific SDSS DR17 optical denominator/proxy diagnostic for packet-gated-paper-to-wiki-reconciliation rp-2. The caption is intentionally narrow: this figure summarizes the cached optical result used for target definition or denominator design, not the unmeasured multi-survey physical claim.}
\label{fig:topic}
\end{figure}

\section{Interpretation and missing observables}\label{sec:missing}
SDSS-only pilot; full proposal requires the additional survey data named in the research-topic page. The full proposal requires: group catalogues, robust central/satellite labels, halo masses, morphology, and multi-redshift selection functions.

Mass and environment are known separable axes in low-redshift galaxy evolution, but a real environmental-quenching analysis requires group/halo and central-satellite information beyond this nearest-neighbour proxy \citep{peng2010,baldry2006,wetzel2013,goubert2024}.


\section{Reproducibility and safety}\label{sec:repro}
This manuscript was generated by local integration run \texttt{INTEGRATED\_9\_PAPERS\_20260709T012051Z}. Inputs are the original RP-1 SDSS query/run directory, the eight-topic SDSS remaining-topic manifest, the overnight shared selection-function packet, the cached-versus-public representativeness packet, Goru robustness outputs, and literature/source placement packets. The output is a local draft PDF and manifest entry only. No public-linked PDF was replaced.

\section{Conclusion}\label{sec:conclusion}
The integration improves the paper package by putting denominator honesty before results. For RP-1, the strongest outcome is a plausible short-paper association draft: broad optical BPT AGN hosts in this capped SDSS emission-line subset have lower catalog sSFR than mass--redshift matched star-forming controls, with robustness caveats. For the other active topics, the correct packaging is a guarded denominator/proxy suite, not eight independent causal feedback papers.


\begin{thebibliography}{}
\bibitem[Abdurro'uf et al.(2022)]{sdssdr17} Abdurro'uf, Accetta, K., Aerts, C., et al. 2022, ApJS, 259, 35
\bibitem[Baldwin et al.(1981)]{baldwin1981} Baldwin, J.~A., Phillips, M.~M., \& Terlevich, R. 1981, PASP, 93, 5
\bibitem[Brinchmann et al.(2004)]{brinchmann2004} Brinchmann, J., Charlot, S., White, S.~D.~M., et al. 2004, MNRAS, 351, 1151
\bibitem[Kauffmann et al.(2003a)]{kauffmann2003bpt} Kauffmann, G., Heckman, T.~M., Tremonti, C., et al. 2003a, MNRAS, 346, 1055
\bibitem[Kauffmann et al.(2003b)]{kauffmann2003mass} Kauffmann, G., Heckman, T.~M., White, S.~D.~M., et al. 2003b, MNRAS, 341, 33
\bibitem[Kewley et al.(2001)]{kewley2001} Kewley, L.~J., Dopita, M.~A., Sutherland, R.~S., Heisler, C.~A., \& Trevena, J. 2001, ApJ, 556, 121
\bibitem[Kewley et al.(2006)]{kewley2006} Kewley, L.~J., Groves, B., Kauffmann, G., \& Heckman, T. 2006, MNRAS, 372, 961
\bibitem[York et al.(2000)]{york2000} York, D.~G., Adelman, J., Anderson, J.~E., Jr., et al. 2000, AJ, 120, 1579

\bibitem[Baldry et al.(2006)]{baldry2006} Baldry, I.~K., Balogh, M.~L., Bower, R.~G., et al. 2006, MNRAS, 373, 469
\bibitem[Goubert et al.(2024)]{goubert2024} Goubert, P.~H., Bluck, A.~F.~L., Piotrowska, J.~M., \& Maiolino, R. 2024, arXiv:2401.12953
\bibitem[Peng et al.(2010)]{peng2010} Peng, Y.-j., Lilly, S.~J., Kovac, K., et al. 2010, ApJ, 721, 193
\bibitem[Wetzel et al.(2013)]{wetzel2013} Wetzel, A.~R., Tinker, J.~L., Conroy, C., \& van den Bosch, F.~C. 2013, MNRAS, 432, 336
\end{thebibliography}

\end{document}
